\section{Objectifs}

À la fin de ce chapitre, tu dois savoir :

\begin{itemize}
\item lire et construire un arbre pondéré ;
\item calculer une probabilité conditionnelle ;
\item utiliser la formule des probabilités totales ;
\item inverser un conditionnement avec la formule de Bayes ;
\item distinguer sensibilité, spécificité et valeurs prédictives ;
\item interpréter une probabilité a priori et une probabilité a posteriori ;
\item raisonner avec des effectifs lorsque les probabilités sont difficiles à interpréter ;
\item expliquer pourquoi une information nouvelle modifie une estimation initiale.
\end{itemize}

\section{1. Situation de départ — un test positif ne signifie pas forcément « malade »}

Un test de dépistage annonce un résultat positif. Est-ce que cela signifie que la personne est presque certainement malade ?

Pas nécessairement.

Le résultat dépend à la fois :

\begin{itemize}
\item de la qualité du test ;
\item de la fréquence de la maladie dans la population ;
\item de la probabilité d'obtenir un faux positif.
\end{itemize}

C'est précisément le type de problème que permet de traiter le raisonnement bayésien.

On note :

\begin{itemize}
\item $M$ : « la personne est malade » ;
\item $\overline M$ : « la personne n'est pas malade » ;
\item $T^+$ : « le test est positif » ;
\item $T^-$ : « le test est négatif ».
\end{itemize}

La question la plus fréquente est alors :

\[
P(M\mid T^+).
\]

Cette probabilité ne doit surtout pas être confondue avec :

\[
P(T^+\mid M).
\]

La première signifie « probabilité d'être malade sachant que le test est positif ».  
La seconde signifie « probabilité d'obtenir un test positif sachant que la personne est malade ».

Ces deux probabilités répondent à deux questions différentes.

\section{2. Probabilité conditionnelle}

Lorsque $P(B)>0$, la probabilité de $A$ sachant $B$ est définie par :

\[
P(A\mid B)=\frac{P(A\cap B)}{P(B)}.
\]

On peut aussi écrire :

\[
P(A\cap B)=P(B)\times P(A\mid B).
\]

Cette deuxième forme est particulièrement utile dans un arbre pondéré : la probabilité d'un chemin est obtenue en multipliant les probabilités portées par ses branches.

\subsection{Interprétation}

Supposons que :

\[
P(M)=0{,}02
\]

et

\[
P(T^+\mid M)=0{,}95.
\]

Alors :

\[
P(M\cap T^+)=P(M)\times P(T^+\mid M)
=0{,}02\times0{,}95
=0{,}019.
\]

Dans l'ensemble de la population, environ $1{,}9\%$ des personnes sont donc à la fois malades et testées positives.

\section{3. L'arbre pondéré}

Un arbre pondéré permet d'organiser les événements successifs.

\coursefigure{/images/mathematiques/terminale-complementaires/inference-bayesienne-arbre.svg}{Arbre pondéré d}{Exemple d'arbre pondéré : prévalence 2 \%, sensibilité 95 \%, spécificité 90 \%.}

Dans cet exemple :

\[
P(M)=0{,}02,
\qquad
P(\overline M)=0{,}98.
\]

La sensibilité vaut $95\%$, donc :

\[
P(T^+\mid M)=0{,}95.
\]

Par complémentarité :

\[
P(T^-\mid M)=1-0{,}95=0{,}05.
\]

La spécificité vaut $90\%$, donc :

\[
P(T^-\mid \overline M)=0{,}90.
\]

Ainsi le taux de faux positifs est :

\[
P(T^+\mid\overline M)=1-0{,}90=0{,}10.
\]

\subsection{Règles de calcul dans l'arbre}

Deux règles sont fondamentales.

\textbf{Règle 1 — sur un chemin, on multiplie.}

Par exemple :

\[
P(M\cap T^+)=P(M)\times P(T^+\mid M).
\]

\textbf{Règle 2 — lorsque plusieurs chemins conduisent au même événement, on additionne.}

L'événement $T^+$ peut arriver :

\begin{itemize}
\item soit parce que la personne est malade et le test est positif ;
\item soit parce qu'elle n'est pas malade mais le test est faussement positif.
\end{itemize}

Donc :

\[
P(T^+)=P(M\cap T^+)+P(\overline M\cap T^+).
\]

\section{4. Formule des probabilités totales}

Si les événements $A_1,\ldots,A_n$ forment une partition de l'univers, alors :

\[
P(B)=\sum_{i=1}^{n} P(A_i)\,P(B\mid A_i).
\]

Dans le cas de la partition $M,\overline M$, on obtient :

\[
P(T^+)
=
P(M)P(T^+\mid M)
+
P(\overline M)P(T^+\mid\overline M).
\]

Avec les valeurs de l'exemple :

\[
P(T^+)
=
0{,}02\times0{,}95
+
0{,}98\times0{,}10.
\]

Ainsi :

\[
P(T^+)=0{,}019+0{,}098=0{,}117.
\]

Environ $11{,}7\%$ de la population obtient donc un test positif.

Attention : cela ne signifie pas que $11{,}7\%$ de la population est malade. Les faux positifs interviennent dans ce résultat.

\section{5. Formule de Bayes}

La formule de Bayes permet d'inverser un conditionnement.

À partir de :

\[
P(A\cap B)=P(A)P(B\mid A)
\]

et de :

\[
P(A\mid B)=\frac{P(A\cap B)}{P(B)},
\]

on obtient :

\[
\boxed{
P(A\mid B)
=
\frac{P(B\mid A)P(A)}{P(B)}
}.
\]

Dans notre contexte :

\[
\boxed{
P(M\mid T^+)
=
\frac{P(T^+\mid M)P(M)}{P(T^+)}
}.
\]

Avec les valeurs calculées précédemment :

\[
P(M\mid T^+)
=
\frac{0{,}95\times0{,}02}{0{,}117}
=
\frac{0{,}019}{0{,}117}
\approx0{,}162.
\]

La probabilité qu'une personne soit malade sachant que son test est positif est donc d'environ :

\[
\boxed{16{,}2\%}.
\]

C'est souvent contre-intuitif : le test détecte pourtant $95\%$ des malades.

La raison est la faible prévalence de la maladie et le nombre relativement important de faux positifs.

\section{6. Exemple développé 1 — raisonner sur 10 000 personnes}

Le raisonnement en effectifs aide à comprendre le résultat précédent.

Considérons une population fictive de 10 000 personnes.

\subsection{Étape 1 — malades et non malades}

La prévalence est $2\%$.

Nombre de malades :

\[
10\,000\times0{,}02=200.
\]

Nombre de non-malades :

\[
10\,000-200=9\,800.
\]

\subsection{Étape 2 — résultats du test chez les malades}

La sensibilité est $95\%$.

Vrais positifs :

\[
200\times0{,}95=190.
\]

Faux négatifs :

\[
200-190=10.
\]

\subsection{Étape 3 — résultats du test chez les non-malades}

La spécificité est $90\%$.

Vrais négatifs :

\[
9\,800\times0{,}90=8\,820.
\]

Faux positifs :

\[
9\,800-8\,820=980.
\]

\subsection{Tableau d'effectifs}

\begin{tabular}{c|c|c|c}
Situation & Test positif & Test négatif & Total \\
Malade & 190 & 10 & 200 \\
Non malade & 980 & 8 820 & 9 800 \\
Total & 1 170 & 8 830 & 10 000 \\
\end{tabular}

Parmi les 1 170 tests positifs, seulement 190 correspondent à des personnes malades.

Donc :

\[
P(M\mid T^+)
=
\frac{190}{1170}
\approx0{,}162.
\]

On retrouve le résultat obtenu avec Bayes.

\subsection{Ce que montre cet exemple}

Le nombre de faux positifs peut dépasser largement le nombre de vrais positifs lorsque l'événement recherché est rare.

La qualité intrinsèque d'un test ne suffit donc pas à déterminer sa valeur prédictive.

\section{7. Sensibilité, spécificité et valeurs prédictives}

Ces quatre notions doivent être clairement distinguées.

\subsection{Sensibilité}

La sensibilité mesure la capacité du test à détecter les personnes malades :

\[
\text{sensibilité}=P(T^+\mid M).
\]

Une sensibilité élevée signifie peu de faux négatifs.

\subsection{Spécificité}

La spécificité mesure la capacité du test à reconnaître les personnes non malades :

\[
\text{spécificité}=P(T^-\mid\overline M).
\]

Une spécificité élevée signifie peu de faux positifs.

\subsection{Valeur prédictive positive}

La valeur prédictive positive est :

\[
\text{VPP}=P(M\mid T^+).
\]

Elle répond à la question : « sachant que le test est positif, quelle est la probabilité d'être malade ? »

\subsection{Valeur prédictive négative}

La valeur prédictive négative est :

\[
\text{VPN}=P(\overline M\mid T^-).
\]

Elle répond à la question : « sachant que le test est négatif, quelle est la probabilité de ne pas être malade ? »

\subsection{Tableau de comparaison}

\begin{tabular}{c|c|c}
Indicateur & Probabilité & Question \\
Sensibilité & $P(T^+\mid M)$ & Le test repère-t-il bien les malades ? \\
Spécificité & $P(T^-\mid\overline M)$ & Le test repère-t-il bien les non-malades ? \\
VPP & $P(M\mid T^+)$ & Un positif est-il réellement malade ? \\
VPN & $P(\overline M\mid T^-)$ & Un négatif est-il réellement non malade ? \\
\end{tabular}

Les deux premières grandeurs caractérisent le comportement du test dans des groupes connus.  
Les deux dernières dépendent fortement de la prévalence.

\section{8. A priori et a posteriori}

En raisonnement bayésien, on distingue deux étapes.

\subsection{Probabilité a priori}

Avant d'observer le résultat du test, la probabilité d'être malade vaut :

\[
P(M)=0{,}02.
\]

Cette probabilité représente l'information disponible avant l'observation.

\subsection{Probabilité a posteriori}

Après un test positif, la probabilité devient :

\[
P(M\mid T^+)\approx0{,}162.
\]

L'observation a donc modifié notre estimation.

On peut résumer le mécanisme ainsi :

\[
\boxed{
\text{a priori}
+
\text{information nouvelle}
\longrightarrow
\text{a posteriori}
}
\]

Bayes ne « crée » pas de certitude. Il met à jour une probabilité à partir d'une information nouvelle.

\section{9. Exemple développé 2 — contrôle qualité en industrie}

Une usine utilise deux machines.

\begin{itemize}
\item la machine $A$ produit $60\%$ des pièces ;
\item la machine $B$ produit $40\%$ des pièces ;
\item $2\%$ des pièces issues de $A$ sont défectueuses ;
\item $5\%$ des pièces issues de $B$ sont défectueuses.
\end{itemize}

On choisit une pièce au hasard et on constate qu'elle est défectueuse.

On cherche :

\[
P(B\mid D).
\]

\subsection{Étape 1 — probabilité d'une pièce défectueuse}

On applique la formule des probabilités totales :

\[
P(D)
=
P(A)P(D\mid A)+P(B)P(D\mid B).
\]

Donc :

\[
P(D)
=
0{,}60\times0{,}02
+
0{,}40\times0{,}05.
\]

Ainsi :

\[
P(D)=0{,}012+0{,}020=0{,}032.
\]

\subsection{Étape 2 — appliquer Bayes}

\[
P(B\mid D)
=
\frac{P(D\mid B)P(B)}{P(D)}.
\]

Donc :

\[
P(B\mid D)
=
\frac{0{,}05\times0{,}40}{0{,}032}
=
\frac{0{,}020}{0{,}032}
=
0{,}625.
\]

Il y a donc $62{,}5\%$ de chances que la pièce défectueuse provienne de la machine $B$.

\subsection{Interprétation}

La machine $B$ ne fabrique que $40\%$ des pièces, mais son taux de défaut est plus élevé. Une pièce défectueuse est donc plus probablement issue de $B$ que de $A$.

\section{10. Effet de la prévalence}

La prévalence joue un rôle majeur dans les valeurs prédictives.

Supposons un test de sensibilité $99\%$ et de spécificité $99\%$.

\subsection{Cas 1 — prévalence 1 \%}

\[
P(M)=0{,}01.
\]

Le taux de faux positifs vaut également $1\%$.

Alors :

\[
P(T^+)
=
0{,}99\times0{,}01
+
0{,}01\times0{,}99
=
0{,}0198.
\]

Donc :

\[
P(M\mid T^+)
=
\frac{0{,}0099}{0{,}0198}
=
0{,}50.
\]

Malgré un test très performant, la valeur prédictive positive n'est que de $50\%$.

\subsection{Cas 2 — prévalence 20 \%}

\[
P(M)=0{,}20.
\]

Alors :

\[
P(T^+)
=
0{,}99\times0{,}20
+
0{,}01\times0{,}80
=
0{,}206.
\]

Et :

\[
P(M\mid T^+)
=
\frac{0{,}198}{0{,}206}
\approx0{,}961.
\]

La VPP est cette fois d'environ $96{,}1\%$.

Le même test peut donc donner des interprétations très différentes selon la population étudiée.

\coursefigure{/images/mathematiques/terminale-complementaires/inference-bayesienne-prevalence-vpp.svg}{Courbe montrant l}{À performances de test constantes, la valeur prédictive positive augmente fortement avec la prévalence.}

Cette courbe permet de visualiser un point essentiel : la qualité du test reste identique, mais l'interprétation d'un résultat positif change avec la fréquence de base de l'événement recherché.

\section{11. Méthode — résoudre un problème de Bayes}

Lorsqu'une question demande une probabilité du type $P(A\mid B)$, utilise la démarche suivante.

\subsection{Étape 1 — définir clairement les événements}

Écris ce que signifient $A$, $B$, $\overline A$, etc.

\subsection{Étape 2 — identifier ce qui est donné}

Repère :

\begin{itemize}
\item les probabilités initiales $P(A)$ ;
\item les probabilités conditionnelles $P(B\mid A)$ ;
\item les complémentaires éventuels.
\end{itemize}

\subsection{Étape 3 — construire un arbre}

L'arbre permet d'éviter les inversions de conditionnement.

\subsection{Étape 4 — calculer la probabilité de l'observation}

Utilise les probabilités totales :

\[
P(B)
=
P(A)P(B\mid A)
+
P(\overline A)P(B\mid\overline A).
\]

\subsection{Étape 5 — appliquer Bayes}

\[
P(A\mid B)
=
\frac{P(B\mid A)P(A)}{P(B)}.
\]

\subsection{Étape 6 — interpréter}

Une réponse comme $0{,}37$ doit être traduite :

« sachant que $B$ est réalisé, la probabilité de $A$ est de $37\%$. »

\subsection{Variante — passer par les effectifs}

Lorsque les pourcentages sont peu intuitifs, choisis une population fictive de 1 000, 10 000 ou 100 000 individus. Le tableau d'effectifs rend souvent le raisonnement beaucoup plus lisible.

\section{12. Indépendance et Bayes}

Deux événements $A$ et $B$ sont indépendants si :

\[
P(A\cap B)=P(A)P(B).
\]

Lorsque $P(B)>0$, cela revient à :

\[
P(A\mid B)=P(A).
\]

Dans ce cas, apprendre que $B$ s'est produit ne modifie pas la probabilité de $A$.

On peut donc voir l'indépendance comme une situation dans laquelle l'information nouvelle $B$ n'apporte aucune information utile sur $A$.

À l'inverse, si :

\[
P(A\mid B)\neq P(A),
\]

alors l'observation de $B$ modifie notre estimation de $A$.

\section{13. Arbre, tableau ou formule ?}

Les trois outils sont complémentaires.

\begin{tabular}{c|c|c}
Outil & Point fort & À privilégier quand… \\
Arbre pondéré & Visualise les conditionnements successifs & les événements sont organisés en étapes \\
Tableau d'effectifs & Rend les proportions concrètes & les probabilités sont contre-intuitives \\
Formule de Bayes & Donne un calcul compact & les probabilités sont déjà identifiées \\
\end{tabular}

Un bon raisonnement peut utiliser les trois successivement.

\section{14. Erreurs fréquentes}

\subsection{Confondre $P(A\mid B)$ et $P(B\mid A)$}

C'est l'erreur principale.

\[
P(A\mid B)
\neq
P(B\mid A)
\]

en général.

Toujours traduire la probabilité en phrase avant de calculer.

\subsection{Oublier les probabilités de base}

Une sensibilité de $99\%$ ne permet pas à elle seule de connaître la probabilité d'être malade après un test positif.

Il faut aussi connaître la prévalence.

\subsection{Oublier un chemin dans les probabilités totales}

Pour calculer $P(T^+)$, il faut compter tous les cas permettant d'obtenir un test positif.

\subsection{Utiliser le complément du mauvais événement}

Si :

\[
P(T^-\mid\overline M)=0{,}90,
\]

alors :

\[
P(T^+\mid\overline M)=0{,}10.
\]

Le complément s'effectue en gardant la même condition.

\subsection{Donner un nombre sans l'interpréter}

Une probabilité conditionnelle est une proportion dans une sous-population. La phrase finale est indispensable.

\section{15. À retenir}

Les quatre relations essentielles sont :

\[
\boxed{
P(A\mid B)=\frac{P(A\cap B)}{P(B)}
}
\]

\[
\boxed{
P(A\cap B)=P(B)P(A\mid B)
}
\]

\[
\boxed{
P(B)=\sum_i P(A_i)P(B\mid A_i)
}
\]

\[
\boxed{
P(A\mid B)=\frac{P(B\mid A)P(A)}{P(B)}
}
\]

Le point central du chapitre est conceptuel :

\[
\boxed{
\text{une probabilité dépend de l'information disponible}
}
\]

La formule de Bayes permet de mettre à jour une probabilité lorsque cette information change.
